\scanpage{352}
\[
 \deleted{Z^j=\bigwedge^j\Psi\times\bigwedge^{\unsure{i+1}}\Phi
 \longrightarrow\bigwedge^{\unsure{*}} C^{i+j}=\bigwedge^{i+j+1}\Phi}
\]
On \unsure{a alors|se donne} \unsure{une loi de|dans les \(Z^i\) la} \(k\)-Algèbre graduée alternée \unsure{des} \(Z^i(C^\bullet)=\bigwedge^i\Psi\), par le produit extérieur \unsure{canonique. Question :|canonique.} quand et comment cela \unsure{se prolonge en une} structure de \(\tfrac12\)-algèbre différentielle ?
\diagram{\begin{tikzcd}[ampersand replacement=\&,row sep=small]
 Z^{i+1}\otimes Z^j \arrow[r,"\text{donné}"] \&Z^{i+j+1}\\
 C^i\otimes Z^j \arrow[u]\arrow[r,"\sigma_{ij}"]\& C^{i+j}\arrow[u]\\
 Z^i\otimes Z^j\arrow[u]\arrow[r,"\text{donné}"']\& Z^{i+j}\arrow[u]\&\unsure{]}\\
 0\arrow[u]\&0\arrow[u]\&
\end{tikzcd}}
L'indétermination pour trouver \(\sigma_{ij}\) est dans
\[
 \operatorname{Hom}_k(Z^{i+1}\otimes Z^j,Z^{i+j})=
\]
\[
 \underset{\displaystyle\simeq\bigwedge^{i+j+1}\Psi_*
 \oplus\,\substack{\text{quelque chose}\\\text{\unsure{d'homotope|d'homotopie} à \(0\)}}}
 {\underbrace{\bigwedge^{i+1}\Psi\otimes\bigwedge^j\Psi}}
 \qquad\bigwedge^{i+j}\Psi
\]
% archive sheet number: 353
